\section{Related Work}

\begin{figure}[t]
  \centering
  \includegraphics[width=\linewidth]{Figures/workflow/workflow2.pdf}
  \caption{\textbf{Workflow.} A short caption is expanded by a VLM into a detailed JSON used by \modelname{} to generate an image. The user can then refine the JSON, and \modelname{} produces a new image that reflects only the requested changes, demonstrating strong disentanglement between modified and preserved elements. 
  }
  \label{fig:workflow}
\end{figure}


\paragraph{Text-to-image models.}
Diffusion models have rapidly become the workhorse of text-to-image synthesis. Early systems, like GLIDE~\cite{nichol2021glide}, Imagen~\cite{saharia2022photorealistic}, and DALL·E~2~\cite{ramesh2022hierarchical}, established the power of conditioning diffusion on strong language encoders. Latent diffusion~\cite{rombach2022high} made large-scale training practical, and SDXL~\cite{podell2023sdxl} pushed UNet-based scaling. A newer wave pivots to transformer backbones and flow-matching objectives: Stable Diffusion 3~\cite{esser2024scaling} adopts DiT-style architectures with bidirectional text–image mixing, while FLUX~\cite{flux2024}, HiDream-I1~\cite{cai2025hidream}, and Qwen-Image~\cite{wu2025qwenimagetechnicalreport} further advance this paradigm. Together, these advances mark a shift toward models that must understand and accurately follow long prompts---the challenge our work is designed to tackle.

\paragraph{Synthetic captions.}

Early text-to-image models were trained on large-scale web datasets such as LAION~\cite{schuhmann2022laion}, where captions derived from HTML alt text provided only coarse and noisy supervision. 
DALL·E~3~\cite{betker2023improving} demonstrated that augmenting training data with synthetically generated, descriptive captions improves both human and automated evaluations, allowing models to capture visual concepts that human annotators often omit. 
Subsequent works~\cite{esser2024scaling, liu2024playground} explored hybrid datasets that mix human-written and synthetic captions to balance linguistic diversity and descriptive precision. 
We extend this idea by using modern vision-language models to produce structured JSON captions that explicitly describe object attributes, spatial relations, composition, and photographic style—capturing fine-grained detail beyond prior open-source efforts. 
The same structured schema is used consistently during training and inference, enabling human-written prompts to be automatically converted into machine-readable structured descriptions.


 

\paragraph{LLM Fusion.} 

Earlier text-to-image models relied on CLIP~\cite{radford2021learning} or T5~\cite{raffel2020exploring} as text encoders~\cite{rombach2022high,podell2023sdxl,saharia2022photorealistic}, mapping raw text to a compact, semantically meaningful representation. Recent work explores leveraging LLMs for text encoding to exploit their richer semantics and compositional reasoning. Because different intermediate LLM layers encode complementary linguistic information~\cite{durrani2020analyzing,dar2022analyzing,rombach2022high}, several methods, e.g., Playground v3~\cite{liu2024playground} and Tang et al.~\cite{tang2025exploring}, fuse multiple layers into diffusion transformers via cross-attention. However, these approaches typically lack bidirectional text–image mixing (i.e., joint attention), which has been shown to improve prompt adherence~\cite{esser2024scaling}. HiDream-I1~\cite{cai2025hidream} use a dual-stream, decoder-only LLM but at high computational cost, since each attention block consumes tokens from both final and intermediate layers. Our proposed DimFusion improves efficiency by leveraging both intermediate and final LLM representations while keeping token length constant.







\paragraph{Cycle-consistency-based evaluation.} 
Cycle consistency is the principle that translating an input from one domain to another and back should approximately reconstruct the original (e.g., image$\rightarrow$text$\rightarrow$image). A prominent example is the self-supervised objective in CycleGAN~\cite{zhu2017unpaired}. In text-to-image generation, recent works use cycle consistency as a reward in reinforcement settings to improve prompt adherence~\cite{meng2025image,bahng2025cycle}, and as an automatic metric for text–image alignment~\cite{huang2025image2text2image}. Meng et~al.~\cite{meng2025image} further propose an automated evaluation based on image regeneration via captions, but without a direct pixel-level or perceptual comparison to the original image. In contrast, we introduce \emph{Text-as-a-Bottleneck Reconstruction (TaBR)}, which given real image, produce a detailed caption, regenerates the image from this caption, and then explicitly compares the reconstruction to the original. This image-grounded protocol provides a more direct measure of expressive power and controllability.


